% =====================================================================
\section{Testing}\label{sec:testing}
All tests are CPU-only (no GPU, no window) and were regenerated for this report by building and running the test binary
(Table~\ref{tab:build}); \texttt{build/test\_results.json} holds the id, description, expected and observed text and the result of each. At the time of
writing \textbf{\nTests{} tests were run, \nPass{} passed and \nFail{} failed}. Table~\ref{tab:testareas} groups them by area; the complete list is Table~\ref{tab:alltests}
(generated from the JSON by \code{report/gen\_tables.py}, never typed by hand).

\begin{table}[h]
\centering\small
\caption{Tests by area (generated).}
\label{tab:testareas}
% generated by report/gen_tables.py -- do not edit
\begin{tabular}{@{}lrr@{}}\toprule
Area (id range) & Tests & Pass\\\midrule
Primitives, transforms, camera maths (T01--T19) & 19 & 19\\
Tracker, occlusion, planner, UAV logic (T20--T32) & 13 & 13\\
Shadows, props, terrain, heat trail, walk camera (T33--T55) & 23 & 23\\
Bunker and terrain (T60--T69) & 10 & 10\\
Base (T70--T79) & 10 & 10\\
Tank and AA gun (T80--T89) & 11 & 11\\
Helicopter and radar (T90--T99) & 10 & 10\\
Combat and game logic (T100--T113) & 14 & 14\\
Application, director, bunker walk (T120--T129) & 10 & 10\\
\midrule
Total & 120 & 120\\
\bottomrule
\end{tabular}

\end{table}

What the tests do \emph{not} cover: anything that needs a GPU (rendering output, GLSL compilation, frame time), the interactive keyboard/mouse path in
\code{App}, the HUD drawing and the video encoder. These parts were checked by looking at captures made with the llvmpipe software renderer (Section~\ref{sec:results}); the old visual checks of v2
are listed in the workbook sheet \emph{Testing} (\code{docs/visual\_tests.json}).
Test T38 was rewritten for v4: the old v3 budget of 1000 leaves no longer applies (the scene has \nLeaves{} leaves) and it now checks that every environment and backdrop leaf is a lab primitive.
Some descriptions still carry v3 wording; for example T33 reports a texel of 2.15\,cm, which corresponds to a 22\,m shadow sphere at $2048^2$, while the renderer now fits a 39\,m sphere (3.8\,cm at $2048^2$). The test
passes but no longer mirrors the renderer (see the audit).

% generated by report/gen_tables.py from build/test_results.json -- do not edit
{\scriptsize\setlength{\tabcolsep}{2.5pt}
\begin{longtable}{@{}l p{0.52\linewidth} p{0.34\linewidth} c@{}}
\caption{All 120 automated tests (\texttt{build/test\_results.json}). The observed value is abridged to 70 characters; the JSON holds the full text.}\label{tab:alltests}\\
\toprule Id & Description & Observed (abridged) & Result\\\midrule
\endfirsthead
\toprule Id & Description & Observed (abridged) & Result\\\midrule
\endhead
\bottomrule\endfoot
T01 & Cube vertex/\allowbreak{}index count & 24 vertices,\allowbreak{} 36 indices,\allowbreak{} 12 triangles & pass\\
T02 & Cube centred at origin with unit size & min (-0.500,\allowbreak{}-0.500,\allowbreak{}-0.500),\allowbreak{} max (0.500,\allowbreak{}0.500,\allowbreak{}0.500) & pass\\
T03 & Cube normals unit length,\allowbreak{} outward,\allowbreak{} CCW winding & 0 wrong-winding,\allowbreak{} 24/\allowbreak{}24 outward,\allowbreak{} max \textbar{}n\textbar{}-1 = 0.00e+00 & pass\\
T04 & Sphere counts follow 2+(S-1)N vertices and 2N(S-1) triangles (N=36,\allowbreak{} S=18) & 614 vertices,\allowbreak{} 1224 triangles & pass\\
T05 & Sphere vertices lie on r = 0.5 and normal = p /\allowbreak{} r & max radius error 2.98e-08,\allowbreak{} max normal error 0.00e+00 & pass\\
T06 & Sphere lower hemisphere is the reflection (y -\textgreater{} -y) of the upper one,\allowbreak{} CCW kept & 614/\allowbreak{}614 vertices mirrored,\allowbreak{} 0 wrong-winding & pass\\
T07 & Cylinder counts 4N+2 vertices,\allowbreak{} 12N indices (N=36) & 146 vertices,\allowbreak{} 432 indices & pass\\
T08 & Cylinder centred at origin,\allowbreak{} axis along Z,\allowbreak{} radius 0.5,\allowbreak{} height 1 & x [-0.500,\allowbreak{}0.500] y [-0.500,\allowbreak{}0.500] z [-0.500,\allowbreak{}0.500] & pass\\
T09 & Cylinder normals unit length and CCW outward winding & 0 wrong-winding,\allowbreak{} max \textbar{}n\textbar{}-1 = 0.00e+00 & pass\\
T10 & Frustum variant (light cone): open,\allowbreak{} side normals tilted,\allowbreak{} CCW & 0 wrong-winding,\allowbreak{} first normal z = -0.425 & pass\\
T11 & All indices valid and no degenerate triangles & indices valid: yes,\allowbreak{} degenerate: 0 & pass\\
T12 & Indexed storage smaller than non-indexed for every primitive & cube 720\textless{}864,\allowbreak{} sphere 29424\textless{}88128,\allowbreak{} cylinder 5232\textless{}10368 B & pass\\
T13 & Terrain disc: radius 36 m,\allowbreak{} 1+R*N vertices,\allowbreak{} N(2R-1) triangles minus the ramp hole,\allowbreak{} all facing +Y & 49921 vertices,\allowbreak{} 97440 triangles (hole 2016),\allowbreak{} radius 36.00,\allowbreak{} 97440/\allowbreak{}97... & pass\\
T14 & Model matrix order M = T*R*S (scale,\allowbreak{} then rotate,\allowbreak{} then translate) & (1.000,\allowbreak{} 4.000,\allowbreak{} 3.000) & pass\\
T15 & Inverse matrix: inverse(M) * M = I & max error 1.19e-07 & pass\\
T16 & Normal matrix keeps normals perpendicular under non-uniform scale & dot = 5.16e-09 & pass\\
T17 & View-projection maps the look-at target to the NDC centre & ndc (0.000,\allowbreak{} 0.000,\allowbreak{} 0.990) & pass\\
T18 & Bunker,\allowbreak{} UAV and Rover are each built from cube + sphere + cylinder & bunker 363 parts,\allowbreak{} UAV 67 parts,\allowbreak{} rover 78 parts & pass\\
T19 & Shading assignment follows the proposal & all leaves match & pass\\
T20 & All world matrices finite after scene construction & all finite & pass\\
T21 & Searchlight is hierarchical: lamp moves with the UAV body & body moved 0.750 m,\allowbreak{} lamp moved 0.776 m & pass\\
T22 & Tracker finds a hot blob and its centroid & found=1,\allowbreak{} area 24,\allowbreak{} u=0.150 v=0.200 & pass\\
T23 & Tracker ignores single-pixel noise (min area 2) & not found & pass\\
T24 & Data association keeps the existing track instead of a bigger distractor & u=-0.82 v=0.77,\allowbreak{} blobs=2 & pass\\
T25 & Segment/\allowbreak{}box occlusion test & through=1 beside=0 & pass\\
T26 & Visibility grid: behind a tall occluder is hidden,\allowbreak{} on the UAV's side visible & occluder found,\allowbreak{} behind (-11.4,\allowbreak{} -5.3) hidden=1,\allowbreak{} front (-6.9,\allowbreak{} -3.2) v... & pass\\
T27 & Flood fill is a valid BFS field and descent reaches a goal (MicroMouse) & goals 213,\allowbreak{} field consistent 1,\allowbreak{} start value 8,\allowbreak{} path 8 cells,\allowbreak{} end val... & pass\\
T28 & Rover drives (LFR controller) from open ground into the UAV's blind spot within 60 s & hidden=1 at (10.10,\allowbreak{} 3.50),\allowbreak{} travelled 7.7 m,\allowbreak{} line under sensors 609/\allowbreak{}... & pass\\
T29 & Pixel-ray ground estimate: image centre of a downward camera hits the point below it & hit=1 at (0.50,\allowbreak{} -0.05,\allowbreak{} 0.50) & pass\\
T30 & 40 s simulation without detections: UAV searches inside the 20 m area,\allowbreak{} above ground & mode SEARCH,\allowbreak{} inside=1,\allowbreak{} finite=1 & pass\\
T31 & Exact occlusion in primitive space (inverse model matrix),\allowbreak{} not just bounding boxes & sphere corner=0 centre=1,\allowbreak{} cylinder corner=0 body=1 & pass\\
T32 & Hysteresis threshold: a big warm blob without a hot core is not a target & blobs=2,\allowbreak{} found=1,\allowbreak{} u=0.58 v=-0.43,\allowbreak{} area 9 & pass\\
T33 & Moon shadow map: the whole 20 m scene sphere projects inside the light's clip volume & max \textbar{}NDC\textbar{} = 1.000,\allowbreak{} texel = 2.15 cm & pass\\
T34 & Shadow cascade snapping: moving the camera shifts the shadow map by whole texels only & shift (3.000,\allowbreak{} -0.000) texels,\allowbreak{} fraction 0.0000 & pass\\
T35 & Searchlight shadow matrix: lamp axis -\textgreater{} map centre,\allowbreak{} 14 deg cone edge inside the map & axis (0.000,\allowbreak{} 0.000),\allowbreak{} edge (0.870,\allowbreak{} 0.000) & pass\\
T36 & Cone (cylinder generator,\allowbreak{} top radius 0): counts,\allowbreak{} centred,\allowbreak{} CCW,\allowbreak{} closed,\allowbreak{} no degenerate triangles & 109 vertices,\allowbreak{} 72 triangles,\allowbreak{} x [-0.500,\allowbreak{}0.500] y [-0.500,\allowbreak{}0.500] z [-0... & pass\\
T37 & Exact cone occlusion in primitive space (pine crowns) & near apex=0,\allowbreak{} axis=1,\allowbreak{} base=1 & pass\\
T38 & New props and backdrop are built only from cube /\allowbreak{} sphere /\allowbreak{} cylinder /\allowbreak{} cone leaves & environment 651 parts \{Cube Sphere Cylinder Cone\},\allowbreak{} backdrop 120 par... & pass\\
T39 & Every material has PBR parameters (roughness,\allowbreak{} metalness in [0,\allowbreak{}1]) and all matrices are finite & all leaves OK & pass\\
T40 & Props sit inside the 20 m disc,\allowbreak{} backdrop parts outside it & farthest prop extent 19.57 m,\allowbreak{} nearest backdrop part 38.01 m & pass\\
T41 & No static ground-level prop is hotter than the rover's motor box or reaches the tracker blob threshold & hottest ground-level part Environment/\allowbreak{}Fire embers = 0.74 & pass\\
T42 & Terrain v2: continuous,\allowbreak{} level bunker pad,\allowbreak{} creek bed,\allowbreak{} track depression,\allowbreak{} berm,\allowbreak{} constant rim & max 1 cm step 0.0055 m,\allowbreak{} pad 0.00000 m,\allowbreak{} creek -0.41 m,\allowbreak{} track -0.03 m... & pass\\
T43 & Terrain mesh density and far-ground clearance & radial 0.277 m,\allowbreak{} rim 0.589 m,\allowbreak{} 49921 vertices,\allowbreak{} 97440 triangles,\allowbreak{} min h... & pass\\
T44 & Bucket-grid line of sight equals the brute-force test over all occluders & 0 mismatches,\allowbreak{} 1656 blocked,\allowbreak{} 516 occluders & pass\\
T45 & 180 s headless run: UAV inside 19.5 m and above ground,\allowbreak{} rover finite inside the free area & UAV max r 15.57 m,\allowbreak{} min clearance 5.62 m; rover max r 16.78 m,\allowbreak{} hidde... & pass\\
T46 & UAV mode logic on the 20 m map: detections -\textgreater{} TRACK,\allowbreak{} loss -\textgreater{} REACQUIRE,\allowbreak{} stays inside & track=1 reacquire=1 inside=1,\allowbreak{} estimate (7.3,\allowbreak{} 2.6) & pass\\
T47 & Heat trail cools exponentially: after tau = 45 s a footprint keeps 1/\allowbreak{}e of its heat & start 0.745,\allowbreak{} after 45 s 0.274,\allowbreak{} ratio 0.368 & pass\\
T48 & Heat trail: deposits are local,\allowbreak{} clamped to 1,\allowbreak{} and hot ground stays below the 0.80 blob threshold & centre 1.00,\allowbreak{} 0.5 m away 0.00,\allowbreak{} max ground heat 0.51 & pass\\
T49 & Scene: after 20 s of driving the rover has left a heat trail on the ground & trail peak 1.00,\allowbreak{} rover at (14.2,\allowbreak{} 8.8) & pass\\
T50 & cameraModeName(Walk) is WALK and Walk comes after Cinematic (existing values unchanged) & "WALK",\allowbreak{} Cinematic = 4,\allowbreak{} Walk = 5 & pass\\
T51 & Walk camera: eye height stays 1.7 +- 0.05 m above the terrain while walking and running & max deviation 0.013 m over 2400 steps & pass\\
T52 & Walk camera cannot fall into the open entrance trench from the side (rails /\allowbreak{} retaining wall) & min x 2.15 m,\allowbreak{} lowest feet -0.08 m & pass\\
T53 & Walk camera is clamped to a 52 m radius & radius 52.00 m,\allowbreak{} eye 1.40 m,\allowbreak{} apron -0.30 m & pass\\
T54 & Walk camera: compute() view sits at the walker (eye 1.7 m up) and mouse-look pitch clamps to +-80 deg & dx 0.000 dz 0.000 dy 0.000,\allowbreak{} pitch 80.0 & pass\\
T55 & C cycles Orbit,\allowbreak{} Free,\allowbreak{} Chase,\allowbreak{} UAV CAM,\allowbreak{} Walk,\allowbreak{} back to Orbit & ORBIT\textgreater{}FREE\textgreater{}CHASE\textgreater{}UAV CAM\textgreater{}WALK\textgreater{}ORBIT & pass\\
T60 & Terrain covers the underground rooms with at least 0.15 m of earth & minimum terrain height -0.217 m & pass\\
T61 & Terrain mesh: 36 m radius,\allowbreak{} hole cut for the entrance trench & 0 inside,\allowbreak{} 1398 around,\allowbreak{} radius 36.00,\allowbreak{} 49921 vertices,\allowbreak{} 97440 triangles & pass\\
T62 & Walk from the ramp top to the command room centre stays on the floors & 367 steps,\allowbreak{} 0 contacts,\allowbreak{} max step 0.023 m,\allowbreak{} final y -3.40 (CIC -3.40),\allowbreak{}... & pass\\
T63 & Collision pushes the walker out of walls and furniture; the surface above is free & wall 1 x=-1.20,\allowbreak{} cross 1,\allowbreak{} console 1,\allowbreak{} roof contact 0 & pass\\
T64 & A closed blast door blocks the walker,\allowbreak{} an open one lets him pass & open z=-2.50,\allowbreak{} closed z=4.30,\allowbreak{} reopened z=-2.50 & pass\\
T65 & Blast door state machine OPEN/\allowbreak{}CLOSING/\allowbreak{}CLOSED/\allowbreak{}OPENING,\allowbreak{} \textasciitilde{}4 s travel,\allowbreak{} reversible & BLAST DOORS OPEN 100\% \textbar{} BLAST DOORS CLOSING 100\% \textbar{} BLAST DOORS CLOS... & pass\\
T66 & Bunker part budget & 363 parts: Blast door assembly 28,\allowbreak{} Bunk room 28,\allowbreak{} Bunker shell 21,\allowbreak{} C... & pass\\
T67 & Bunker parts are lab primitives with name and purpose,\allowbreak{} Flat shading,\allowbreak{} occluders only above ground & 0 foreign,\allowbreak{} 0 undocumented,\allowbreak{} 363/\allowbreak{}363 flat,\allowbreak{} 11 occluders (0 buried) & pass\\
T68 & Rig interface: fans,\allowbreak{} beacon above ground,\allowbreak{} ramp cube matches the walk floor,\allowbreak{} door camera & fans 1,\allowbreak{} beacon y 4.30,\allowbreak{} ramp error 0.000 m,\allowbreak{} eye (0.0 -0.3 9.0) & pass\\
T69 & Walk space inventory & 12 floors,\allowbreak{} 36 walls,\allowbreak{} 5 indoor volumes,\allowbreak{} outside height 0.091 (terrai... & pass\\
T70 & every Base part is a lab primitive with a name and purpose & 1250 parts,\allowbreak{} 0 violations & pass\\
T71 & no Base part intersects the helipad,\allowbreak{} radar,\allowbreak{} AA pit,\allowbreak{} tank spot,\allowbreak{} road or trench strip & 0 hits  & pass\\
T72 & perimeter wall is closed except one gate opening centred on +Z (90 deg),\allowbreak{} 9..11 m wide & 204 cells; largest gap 11.0 m at 90.8 deg,\allowbreak{} next 1.00 m & pass\\
T73 & boom barrier: 0 to 85 deg in 2.5 s with smoothstep easing,\allowbreak{} arm box only when lowered,\allowbreak{} input toggles & travel 2.50 s,\allowbreak{} half-way at 1.25 s,\allowbreak{} final 85.0 deg,\allowbreak{} tip y 6.04,\allowbreak{} clos... & pass\\
T74 & waving flag: travelling wave,\allowbreak{} strip rotations bounded,\allowbreak{} flag never stretches or collapses & 12 strips,\allowbreak{} max strip yaw 6.0 deg,\allowbreak{} pitch 1.5 deg,\allowbreak{} tip distance 2.50.... & pass\\
T75 & windsock: 5 hierarchical cone segments that droop more in weak wind & 5 chained segments,\allowbreak{} segment pitch 16.0 deg at wind 0.36 vs 0.7 deg... & pass\\
T76 & Base part counts (static apron + controllables) & static 1168 + controllable 82 = 1250,\allowbreak{} 3 controllables & pass\\
T77 & Base parts are never occluders and stay on the terrain mesh & occluders 0,\allowbreak{} parts beyond r = 36: 0,\allowbreak{} farthest centre 35.5 m & pass\\
T78 & compound: trench strip and bunker footprint free of props; \textgreater{}= 70\% of the v3 occluders kept & strip 0,\allowbreak{} footprint 0 ,\allowbreak{} occluders 505 (v3: 477) & pass\\
T79 & backdrop starts beyond the terrain mesh (r \textgreater{}= 36) and its ground stays just below the apron level & nearest 38.0 m,\allowbreak{} far ground top -0.35,\allowbreak{} solid flags 0 & pass\\
T80 & all 80 track links stay on the track path while driving & max distance 0.00000 m,\allowbreak{} parts: 226 (tracks 80) & pass\\
T81 & track belt advances by speed*dt (link 0 sits at its arc-length position) & advance 16.0000 m vs 16.0000 m (v=8.00 m/\allowbreak{}s,\allowbreak{} perimeter 14.342 m),\allowbreak{} li... & pass\\
T82 & skid steer: equal track speeds drive a straight line & yaw change 0.0000 deg,\allowbreak{} path bearing 30.00 deg (heading 30),\allowbreak{} distanc... & pass\\
T83 & skid steer: opposite track speeds pivot with yaw rate (vr - vl) /\allowbreak{} track width & -62.50 deg/\allowbreak{}s vs -62.50 deg/\allowbreak{}s (vl 1.20,\allowbreak{} vr -1.20),\allowbreak{} drift 0.120 m & pass\\
T84 & turret slews 40 deg/\allowbreak{}s,\allowbreak{} gun 20 deg/\allowbreak{}s,\allowbreak{} elevation limited to -8..+20 deg & yaw 40.0 deg/\allowbreak{}s,\allowbreak{} pitch after 1 s 20.0,\allowbreak{} max 20.00,\allowbreak{} min -8.00,\allowbreak{} huge mo... & pass\\
T84b & turret stays stabilised when the hull turns while possessed & hull yaw -114.8 deg,\allowbreak{} turret world yaw -40.000 -\textgreater{} -40.000 & pass\\
T85 & main gun: 120 m/\allowbreak{}s shell,\allowbreak{} 0.35 m recoil returning in 0.6 s,\allowbreak{} 3.5 s reload,\allowbreak{} 18 rounds & shells 1 at 120.0 m/\allowbreak{}s,\allowbreak{} recoil peak 0.350,\allowbreak{} after 0.0000,\allowbreak{} rounds afte... & pass\\
T86 & hull height follows the terrain (mean of the 4 track corners) & worst error 0.0023 m (terrain relief under the hull up to 0.41 m) & pass\\
T87 & collision against an obstacle box stops the tank; perimeter limit r \textless{}= 31 m & front edge z 7.999 m,\allowbreak{} speed 0.00 m/\allowbreak{}s; perimeter run r = 31.000 m & pass\\
T88 & AA gun: 0.5 s spin-up,\allowbreak{} then 30 rounds/\allowbreak{}s,\allowbreak{} spin coasts down in 1.5 s & first round at 0.517 s,\allowbreak{} 60 rounds in 2 s,\allowbreak{} spin 5.00 -\textgreater{} 2.50 -\textgreater{} 0.00... & pass\\
T89 & AA gun: elevation -5..+85 deg,\allowbreak{} 60 deg/\allowbreak{}s azimuth,\allowbreak{} idle sweep never fires & pitch max 85.00 min -5.00,\allowbreak{} azimuth 60.0 deg/\allowbreak{}s,\allowbreak{} idle: moved 157.4 de... & pass\\
T90 & helicopter rotor spools up in \textasciitilde{}6 s after the start command & rest 0\%,\allowbreak{} startup 1,\allowbreak{} ready after 6.00 s,\allowbreak{} rotor 100\% & pass\\
T91 & 5 blades evenly spaced (72 deg) and the angle integrates rpm * dt & spool 5445,\allowbreak{} 2 s 3600.0,\allowbreak{} spacing error 0.000,\allowbreak{} blade1 45.0 vs 45.0 & pass\\
T92 & attitude follows acceleration but stays bounded; speed capped at 25 m/\allowbreak{}s & max pitch 22.2,\allowbreak{} roll 25.5,\allowbreak{} nose-down 21.7,\allowbreak{} speed 24.5 m/\allowbreak{}s & pass\\
T93 & altitude limited to 0.5..40 m AGL while airborne,\allowbreak{} radius \textless{}= 60 m & max AGL 40.00,\allowbreak{} min AGL 0.50,\allowbreak{} max r 60.00 & pass\\
T94 & landing only from low speed /\allowbreak{} height on flat ground; touches down and spools down & high 1,\allowbreak{} fast 1,\allowbreak{} accepted 1,\allowbreak{} parked 1,\allowbreak{} rotor 0\%,\allowbreak{} steepest slope 0.31... & pass\\
T95 & not possessed: parked helicopter stays still; released in flight it hovers in place & parked still 1; flying: speed 0.00,\allowbreak{} alt 14.5 -\textgreater{} 14.5,\allowbreak{} drift 0.019,\allowbreak{}... & pass\\
T96 & helicopter: lab primitives only,\allowbreak{} every node named,\allowbreak{} blur disc fades in,\allowbreak{} parked on the pad with an obstacle box & leaves 68,\allowbreak{} lab 1,\allowbreak{} named 1,\allowbreak{} blur 0.00 -\textgreater{} 0.26,\allowbreak{} boxes 2,\allowbreak{} pos (20.3,\allowbreak{} 1... & pass\\
T97 & radar ROTATE mode turns 12 rpm (360 deg in 5 s) & swept 360.00 deg in 5 s,\allowbreak{} mode ROTATE & pass\\
T98 & radar SECTOR mode sweeps +-45 deg around the aimed centre & mode SECTOR,\allowbreak{} centre 165.6,\allowbreak{} range -45.0 .. 45.0 & pass\\
T99 & radar tilt limited to -5..+20 deg; antenna = 7 x 3 curved tiles; lab primitives only & tilt -5.0 .. 20.0,\allowbreak{} tiles 21,\allowbreak{} leaves 102,\allowbreak{} lab 1,\allowbreak{} named 1 & pass\\
T100 & Homing round launched 30 deg off the line of sight detonates within 0.6 m of the laser spot & blast yes at 0.32 s,\allowbreak{} 0.50 m from the spot & pass\\
T101 & Tank shell is ballistic: range on level ground matches v\textasciicircum{}2 sin(2a)/\allowbreak{}g within 2\% & 203.8 m (-0.24\%) & pass\\
T102 & A tracer cannot tunnel through a small target (segment vs sphere per step) & direct hit id 7 & pass\\
T103 & Explosion: fire,\allowbreak{} smoke,\allowbreak{} sparks,\allowbreak{} flash,\allowbreak{} shock wave,\allowbreak{} dust,\allowbreak{} embers + decaying light + scorch & 7 types,\allowbreak{} light 1 -\textgreater{} 0,\allowbreak{} 1 scorch,\allowbreak{} shake 0.81 & pass\\
T104 & Particle count never exceeds the cap (oldest smoke dropped first) & 4000 & pass\\
T105 & Blast damage radius: destroyed at 0.5 r,\allowbreak{} untouched at 3 r & intruder 0 alive=0,\allowbreak{} intruder 1 alive=1,\allowbreak{} kills 1 & pass\\
T106 & Debris: rigid pieces thrown up,\allowbreak{} all at rest on (not under) the terrain within 8 s & 23 pieces,\allowbreak{} peak 3.25 m,\allowbreak{} 23 resting,\allowbreak{} 23 above ground & pass\\
T107 & Wreck life cycle: burning wreck -\textgreater{} removed after 25 s -\textgreater{} respawn 3 s later & wreck=1 gone=1 alive=1 & pass\\
T108 & Launcher: no round before the target has been tracked for 1.5 s & first request at 1.53 s (mode TRACK) & pass\\
T109 & Target discrimination: a blob on a known wreck is ignored (no track,\allowbreak{} no launch) & mode SEARCH,\allowbreak{} lock 0.00 s,\allowbreak{} request 0 & pass\\
T110 & Friend-or-foe: blobs within 4.5 m of a friendly vehicle are not targets & mode SEARCH,\allowbreak{} request 0 & pass\\
T111 & Keep-out: no launch at a target on the bunker's own entrance ramp & mode TRACK,\allowbreak{} any request 0 & pass\\
T112 & UAV launcher: 4 rail rounds,\allowbreak{} hidden at launch,\allowbreak{} replaced after the 10 s reload & launched 4,\allowbreak{} ready 0,\allowbreak{} hidden 1; after reload ready 4 visible 1 & pass\\
T113 & Full base simulated 120 s: finite state,\allowbreak{} \textgreater{}= 7 possessable objects (UAV + \textgreater{}= 6 new) & 9 controllables (7 possessable),\allowbreak{} 2983 leaves,\allowbreak{} finite 1,\allowbreak{} kills 0,\allowbreak{} ro... & pass\\
T120 & Walk camera follows the entrance ramp floor going down (feet on the slope,\allowbreak{} eye 1.7 m above) & max feet deviation 0.000 m,\allowbreak{} max eye deviation 0.013 m,\allowbreak{} z = 6.39 m,\allowbreak{}... & pass\\
T121 & Walk camera cannot pass a closed blast door box and stays inside the ramp walls & stopped at z = 4.400 m,\allowbreak{} max \textbar{}x\textbar{} 1.200 m & pass\\
T122 & With the doors open the walker goes down the ramp,\allowbreak{} through the door and reaches the CIC & z = -6.01 m,\allowbreak{} feet -3.40 m,\allowbreak{} eye -1.70 m,\allowbreak{} smallest ceiling gap 0.90 m & pass\\
T123 & Walk camera is clamped to 52 m from the origin and the eye to 0.15 m under a low ceiling & radius 52.00 m,\allowbreak{} eye 6.350 m & pass\\
T124 & Camera shake: none at amount 0,\allowbreak{} small and repeatable at amount 1 & rest 0.00000,\allowbreak{} shaken 0.0150,\allowbreak{} repeat 0.000000 & pass\\
T125 & Director: \textasciitilde{}120 s,\allowbreak{} 12 shots,\allowbreak{} contiguous (every moment belongs to exactly one shot,\allowbreak{} fade dips at each cut) & length 120.00 s,\allowbreak{} 12 shots,\allowbreak{} sum 120.00 s,\allowbreak{} contiguous 1,\allowbreak{} finite 1,\allowbreak{} ab... & pass\\
T126 & Director bunker shot: walk path descends the ramp,\allowbreak{} passes the door plane and ends in the rooms & z end -6.00 m,\allowbreak{} eye end -1.71 m,\allowbreak{} lowest eye -1.72 m,\allowbreak{} max vertical st... & pass\\
T127 & Possession keys 5 6 7 8 9 0 map to tank,\allowbreak{} helicopter,\allowbreak{} AA gun,\allowbreak{} radar,\allowbreak{} blast doors,\allowbreak{} UAV (found by name) & keys 1,\allowbreak{} names 1,\allowbreak{} fallback 1 & pass\\
T128 & [ ] cycle all possessable controllables (flag skipped,\allowbreak{} wraps); walk interaction picks the door /\allowbreak{} barrier within 3.5 m & first 1 last 7 wrap 1,\allowbreak{} door 5,\allowbreak{} barrier 6,\allowbreak{} far -1 & pass\\
T129 & Director scripts: doors opened early,\allowbreak{} helicopter lifts,\allowbreak{} tank drives,\allowbreak{} AA fires,\allowbreak{} UAV strike at 81 s,\allowbreak{} nobody during the walk & slots 4 1 0 2 -1,\allowbreak{} lift 1.0,\allowbreak{} forward 1.0,\allowbreak{} fire 1,\allowbreak{} strike 1,\allowbreak{} doors ac... & pass\\
\end{longtable}}

